\chapter{Conclusion}
\label{chap:conc}

This project set out to demonstrate that a carefully implemented 2D Attention U-Net, trained
and evaluated under a rigorous per-volume protocol, can achieve competitive multi-organ
abdominal CT segmentation results using only the compute available on a free-tier cloud GPU.
The final model (Run~017) achieves a mean Dice of 0.8769 and an HD95 of 7.59~mm across 13
abdominal organs on the FLARE22 benchmark --- and does so in under 80 minutes of training
time on a single T4 GPU.

Several design decisions contributed to this outcome. Attention gates on all skip connections
improved performance on small organs by suppressing irrelevant background activations in the
decoder. The combined cross-entropy and Dice loss provided stable gradients from the start of
training while counteracting the severe class imbalance inherent to abdominal CT. Cosine
annealing allowed the learning rate to decay smoothly, enabling continued fine-tuning through
the full 60 epochs without overshooting. Test-Time Augmentation with a four-way flip ensemble
added approximately $+$0.06 Dice at zero additional training cost. And the corrected per-volume
evaluation protocol --- replacing the na\"\i{}ve per-slice average that inflated scores through
trivially empty slices --- ensured that every number in this report reflects genuine
segmentation quality.

The comparison with nnU-Net 3D is instructive. On large, well-defined organs the 2D model
comes within 0.01 Dice of the 3D baseline. On 7 of 13 organs it numerically exceeds the published nnU-Net~3D baseline --- a result we
attribute primarily to the 2D model's ability to process full-resolution in-plane slices
without the patch-level downsampling that 3D methods require to fit GPU memory. The
comparison is indicative rather than controlled, as the nnU-Net figures come from a different
evaluation protocol. The pancreas and duodenum remain the
harder cases: their irregular shapes and poor tissue contrast demand the inter-slice context
that only volumetric models can provide natively.

Three broader lessons emerge. First, evaluation rigour matters as much as architecture: fixing
the per-slice Dice bug shifted the reported score from an inflated $\sim$0.88 to a meaningful
0.8164 (no TTA), a correction that is honest but uncomfortable. Second, data volume is the
dominant factor at this scale --- scaling from 30 to 50 training volumes added more Dice than
any single architectural or scheduling change. Third, the gap between 2D and 3D methods is not
fixed; it narrows considerably when 2D methods are given sufficient data, correct evaluation,
and modern training recipes.

In summary, the Attention U-Net trained in this project is a practical, accurate, and efficient
solution for multi-organ abdominal CT segmentation under real hardware constraints.
